\scanpage{147}
% archive sheet number: 146
% L1–6
d'orientation \deleted{\unsure{$\varepsilon_i$}} $\mathrm{sg}_i$ $(i\in\Pi)$ sur $G_0$ sont les $=$s (cf plus haut). Si $n$ impair, \unsure{cel|cela} provient du fait que sur $G_0$, engendré par des \deleted{\unsure{rel.}}\unsure{gén.} $\tau_0,\tau_1,\tau_\infty$ satisfaisant aux relations (40)(41), il y a \deleted{sous-groupe}\inserted{alors un} unique sous-groupe d'indice 2, de façon plus précise
% L7
\[
 (G_0)_{\rm ab}\simeq \deleted{\unsure{Z}}\pm1 \qquad\text{si }n\text{ impair}
 \tag*{(49)}
\]
% L8–12
$[\![$\uline{NB} \unsure{En tous cas|En tout cas}, $(G_0)_{\rm ab}$ est \deleted{\unsure{eng.}} le groupe abélien engendré par les trois él. d'ordre 2, $\dot\tau_0,\dot\tau_1,\dot\tau_\infty$ \inserted{donc est un quotient de $\mathbb F_2^3$} et $\dot\rho_\infty^3=1$ \deleted{\unsure{se tr}} implique $\dot\rho_\infty=1$ i.e. $\dot\tau_1\dot\tau_0=1$ i.e. $\dot\tau_0=\dot\tau_1$, donc $G_{0\,\rm ab}$ est engendré par $\dot\tau_1,\dot\tau_\infty$, avec la relation
% L13
\[
(\dot\tau_1\dot\tau_\infty)^n=1,\quad\text{i.e. }\dot\tau_1=\dot\tau_\infty\quad\text{si }n\text{ impair}
\tag*{(50)}
\]
% L14–15
ce qui montre notre assertion pour $n$ impair. Pour $n$ pair, on aura
% L16
\[
(G_0)_{\rm ab}\simeq\{\dot\tau_1,\dot\tau_\infty\mid\dot\tau_1^2=\dot\tau_\infty^2=\dot\tau_1\dot\tau_\infty\}\simeq\mathbb F_2\times\mathbb F_2,\tag*{(51)}
\]
% L17–19
il y a donc \inserted{un ens $X$ de} \uline{trois} caractères \deleted{d'ordre} non triviaux $G_0\longrightarrow\pm1$, qui a priori peuvent être permutés entre eux par \inserted{$G$, \deleted{via $G/G_0$} opérant sur $G_{0\,\rm ab}$ via $G/G_0\simeq$}\deleted{\unsure{Z}}\,$Z'_n$. Mais comme $Z\subset G$ centralise $G_0$, $\operatorname{Im}(\underset{\simeq Z'_n}{Z}\mathrel{-}Z'_n)$ (par $\lambda\longmapsto\lambda^2$) opère trivialement, donc $Z'_n$ opère sur $X$ via
% L22–25
$Z'_n/{Z'_n}^2$ \deleted{\unsure{d'ordre}}, dont l'image dans $\unsure{\mathfrak S_X|\mathscr S_X}\simeq\mathfrak S_3$ est donc un 2-groupe \deleted{\unsure{et} on voudrait prouver}\inserted{donc ou bien réduit} à l'élément neutre (et alors on gagne), ou bien une transposition \deleted{et on voudrait que}\inserted{Dans ce dernier cas}, on reste
